% !TEX root = ../main.tex
\section{Quantizing the agent-to-server transmissions}
\label{sec:quantization}

FedNPG-ADMM reduces communication by shrinking \emph{what} is transmitted from
a matrix to two vectors. A second axis, which the original work does not
explore, is to shrink \emph{how precisely} those vectors are transmitted. The
paper reports overhead at 64 bits per parameter. If performance survives at a
handful of bits per parameter, the saving multiplies the $\Oh(d^2) \to \Oh(d)$
one.

This is the novel contribution of the project. The precedent is Q-GADMM
\citep{elgabli2021qgadmm}, which \citet{lan2023fednpgadmm} cite themselves, and
which quantizes the vectors an ADMM scheme transmits in the \emph{supervised}
setting. We carry that idea into the reinforcement learning setting, where the
transmitted vectors are gradient and consensus directions computed from a
non-stationary sampling distribution. We found no prior work quantizing
agent-to-server transmissions for federated policy gradient with
continuous-control experiments. The nearest results are for compressed temporal
difference learning with linear function approximation
\citep{mitra2024compressedtd}, which is theory-only and not policy gradient.

\subsection{Design}

A quantizer is a \texttt{Transform}: \texttt{tensor} $\rightarrow$
(\texttt{tensor}, \texttt{bits\_per\_element}), attached to the
\texttt{Channel} through a factory that \texttt{train()} accepts. It quantizes
and dequantizes in one call and declares the bit width the wire would have
needed. No file in the algorithm proper knows quantization exists.

Two properties make the experiment measure something. The degraded tensor is
what the server aggregates, because downstream code uses the return value of
\texttt{send} and never the original, so the quantization error enters the
update to $\bthe$. And the declared \texttt{bits\_per\_element} is what the
byte meter uses, so the cost axis moves with the fidelity. Per iteration, with
$R$ ADMM rounds, $N$ agents and $b$ bits, the uplink is
\begin{equation}
  \text{bytes} \;=\; N \,(1 + R)\, d \, b / 8,
  \label{eq:quantbytes}
\end{equation}
since $\bg_i$ is sent once per iteration and $\by_i$ once per round. We use
$R = 1$ throughout, so each agent uplinks $2d$ values per iteration and the
saving relative to the paper's 64-bit convention is exactly $64/b$: $8\times$
at 8 bits and $10.7\times$ at 6. Measured byte counts match
\eqref{eq:quantbytes} exactly.

\paragraph{Stochastic uniform quantization.} For a tensor $x$ at $b$ bits, with
$L = 2^b - 1$ levels and step $s = (\max x - \min x)/L$, the value is placed at
$\lfloor u \rfloor + \mathrm{Bernoulli}(u - \lfloor u \rfloor)$ where
$u = (x - \min x)/s$. Rounding stochastically instead of to nearest makes the
estimator unbiased, $\Exp[\hat{x}] = x$, with per-element variance at most
$s^2/4$. That property is the reason to prefer it here: a deterministic
rounding error is systematic and would accumulate across iterations, whereas an
unbiased one leaves the server receiving the correct gradient and consensus
direction in expectation. The quantizer draws from a dedicated
\texttt{torch.Generator} rather than the global stream, so runs at different
bit widths see identical action-sampling sequences and differ only in the
quantization.

\paragraph{Protocol.} We run the full grid $N \in \{2, 4, 8\}$ $\times$
$b \in \{6, 8\}$ on all three environments, with the same five seeds and the
same hyperparameters as Section~\ref{sec:reproduction} (90 runs). The
unquantized reference is therefore the matching column of
Table~\ref{tab:reproduction}, and every quantized cell differs from it only in
the uplink channel.

\subsection{Results}

\begin{table}[t]
  \centering
  \caption{Final reward under stochastic quantization, mean $\pm$ standard
  deviation across five seeds, computed as in Table~\ref{tab:reproduction}. The
  64-bit rows are the unquantized $N = 2, 4, 8$ columns of
  Table~\ref{tab:reproduction}. Uplink is per agent per iteration, $2db/8$
  bytes, and does not depend on $N$; the federation total is $N$ times it,
  \eqref{eq:quantbytes}.}
  \label{tab:quant}
  \begin{tabular}{llrrrrr}
    \toprule
    Environment & Bits & $N=2$ & $N=4$ & $N=8$ & Uplink (bytes) & Saving \\
    \midrule
    \multirow{3}{*}{Swimmer-v4}
      & 64 & $126.0 \pm 1.3$ & $127.0 \pm 0.8$ & $128.4 \pm 0.8$ & \num{77888} & $1\times$ \\
      & 8  & $126.1 \pm 0.8$ & $127.1 \pm 0.5$ & $128.6 \pm 0.4$ & \num{9736}  & $\mathbf{8.0\times}$ \\
      & 6  & $128.9 \pm 7.4$ & $125.8 \pm 3.4$ & $127.8 \pm 0.7$ & \num{7302}  & $\mathbf{10.7\times}$ \\
    \midrule
    \multirow{3}{*}{Hopper-v4}
      & 64 & $2140 \pm 386$ & $1658 \pm 273$ & $1605 \pm 159$ & \num{295008} & $1\times$ \\
      & 8  & $1901 \pm 593$ & $1755 \pm 302$ & $1713 \pm 150$ & \num{36876}  & $\mathbf{8.0\times}$ \\
      & 6  & $1573 \pm 817$ & $1853 \pm 551$ & $1837 \pm 248$ & \num{27657}  & $\mathbf{10.7\times}$ \\
    \midrule
    \multirow{3}{*}{Humanoid-v4}
      & 64 & $3391 \pm 796$ & $5343 \pm 121$ & $6329 \pm 334$ & \num{11633184} & $1\times$ \\
      & 8  & $2956 \pm 593$ & $5244 \pm 237$ & $5997 \pm 259$ & \num{1454148}  & $\mathbf{8.0\times}$ \\
      & 6  & $497 \pm 34$   & $676 \pm 73$   & $1514 \pm 526$ & \num{1090611}  & $10.7\times$ \\
    \bottomrule
  \end{tabular}
\end{table}

\begin{figure}[t]
  \centering
  \begin{subfigure}[b]{0.49\textwidth}
    %% Return against cumulative uplink traffic at N=8, Swimmer-v4, per bit width.
%% Mean over five seeds; band spans min to max. The 64-bit curve is the N=8
%% series of Figure 2, so the horizontal gap between curves is the saving.
%% Data: figures/quant_comm_swimmer.dat and figures/communication_swimmer.dat
%% (scripts/export_figures.py).
\begin{tikzpicture}
\begin{axis}[
    width=\linewidth,
    height=0.78\linewidth,
    xlabel={Cumulative uplink (bytes)},
    ylabel={Average batch return},
    xmode=log,
    ymin=0, ymax=140,
    ytick distance=35,
    ymajorgrids=true, xmajorgrids=true, grid style=dashed,
    tick label style={font=\footnotesize},
    label style={font=\footnotesize},
    legend style={legend columns=1, font=\footnotesize, cells={anchor=west},
                  anchor=north west, at={(0.02,0.98)}},
]

\addplot[name path=N8qlo, draw=none, forget plot]
    table[x=N8_bytes, y=N8_lo] {figures/communication_swimmer.dat};
\addplot[name path=N8qhi, draw=none, forget plot]
    table[x=N8_bytes, y=N8_hi] {figures/communication_swimmer.dat};
\addplot[colorD, opacity=0.15, forget plot] fill between[of=N8qlo and N8qhi];
\addplot[colorD, line width=0.9pt, solid]
    table[x=N8_bytes, y=N8_mean] {figures/communication_swimmer.dat};
\addlegendentry{64 bit};

\addplot[name path=q8qlo, draw=none, forget plot]
    table[x=q8_bytes, y=q8_lo] {figures/quant_comm_swimmer.dat};
\addplot[name path=q8qhi, draw=none, forget plot]
    table[x=q8_bytes, y=q8_hi] {figures/quant_comm_swimmer.dat};
\addplot[colorA, opacity=0.15, forget plot] fill between[of=q8qlo and q8qhi];
\addplot[colorA, line width=0.9pt, dashed]
    table[x=q8_bytes, y=q8_mean] {figures/quant_comm_swimmer.dat};
\addlegendentry{8 bit};

\addplot[name path=q6qlo, draw=none, forget plot]
    table[x=q6_bytes, y=q6_lo] {figures/quant_comm_swimmer.dat};
\addplot[name path=q6qhi, draw=none, forget plot]
    table[x=q6_bytes, y=q6_hi] {figures/quant_comm_swimmer.dat};
\addplot[colorB, opacity=0.15, forget plot] fill between[of=q6qlo and q6qhi];
\addplot[colorB, line width=0.9pt, dash dot]
    table[x=q6_bytes, y=q6_mean] {figures/quant_comm_swimmer.dat};
\addlegendentry{6 bit};
\end{axis}
\end{tikzpicture}

    \caption{Swimmer-v4}
  \end{subfigure}
  \hfill
  \begin{subfigure}[b]{0.49\textwidth}
    %% Return against cumulative uplink traffic at N=8, Hopper-v4, per bit width.
%% Mean over five seeds; band spans min to max. The 64-bit curve is the N=8
%% series of Figure 2, so the horizontal gap between curves is the saving.
%% Data: figures/quant_comm_hopper.dat and figures/communication_hopper.dat
%% (scripts/export_figures.py).
\begin{tikzpicture}
\begin{axis}[
    width=\linewidth,
    height=0.78\linewidth,
    xlabel={Cumulative uplink (bytes)},
    ylabel={Average batch return},
    xmode=log,
    ymin=0, ymax=3600,
    ytick distance=900,
    scaled y ticks=base 10:-3,
    ymajorgrids=true, xmajorgrids=true, grid style=dashed,
    tick label style={font=\footnotesize},
    label style={font=\footnotesize},
    legend style={legend columns=1, font=\footnotesize, cells={anchor=west},
                  anchor=north west, at={(0.02,0.98)}},
]

\addplot[name path=N8qlo, draw=none, forget plot]
    table[x=N8_bytes, y=N8_lo] {figures/communication_hopper.dat};
\addplot[name path=N8qhi, draw=none, forget plot]
    table[x=N8_bytes, y=N8_hi] {figures/communication_hopper.dat};
\addplot[colorD, opacity=0.15, forget plot] fill between[of=N8qlo and N8qhi];
\addplot[colorD, line width=0.9pt, solid]
    table[x=N8_bytes, y=N8_mean] {figures/communication_hopper.dat};
\addlegendentry{64 bit};

\addplot[name path=q8qlo, draw=none, forget plot]
    table[x=q8_bytes, y=q8_lo] {figures/quant_comm_hopper.dat};
\addplot[name path=q8qhi, draw=none, forget plot]
    table[x=q8_bytes, y=q8_hi] {figures/quant_comm_hopper.dat};
\addplot[colorA, opacity=0.15, forget plot] fill between[of=q8qlo and q8qhi];
\addplot[colorA, line width=0.9pt, dashed]
    table[x=q8_bytes, y=q8_mean] {figures/quant_comm_hopper.dat};
\addlegendentry{8 bit};

\addplot[name path=q6qlo, draw=none, forget plot]
    table[x=q6_bytes, y=q6_lo] {figures/quant_comm_hopper.dat};
\addplot[name path=q6qhi, draw=none, forget plot]
    table[x=q6_bytes, y=q6_hi] {figures/quant_comm_hopper.dat};
\addplot[colorB, opacity=0.15, forget plot] fill between[of=q6qlo and q6qhi];
\addplot[colorB, line width=0.9pt, dash dot]
    table[x=q6_bytes, y=q6_mean] {figures/quant_comm_hopper.dat};
\addlegendentry{6 bit};
\end{axis}
\end{tikzpicture}

    \caption{Hopper-v4}
  \end{subfigure}

  \vspace{0.6em}
  \begin{subfigure}[b]{0.49\textwidth}
    %% Return against cumulative uplink traffic at N=8, Humanoid-v4, per bit width.
%% Mean over five seeds; band spans min to max. The 64-bit curve is the N=8
%% series of Figure 2, so the horizontal gap between curves is the saving.
%% Data: figures/quant_comm_humanoid.dat and figures/communication_humanoid.dat
%% (scripts/export_figures.py).
\begin{tikzpicture}
\begin{axis}[
    width=\linewidth,
    height=0.78\linewidth,
    xlabel={Cumulative uplink (bytes)},
    ylabel={Average batch return},
    xmode=log,
    ymin=0, ymax=7000,
    ytick distance=1750,
    scaled y ticks=base 10:-3,
    ymajorgrids=true, xmajorgrids=true, grid style=dashed,
    tick label style={font=\footnotesize},
    label style={font=\footnotesize},
    legend style={legend columns=1, font=\footnotesize, cells={anchor=west},
                  anchor=north west, at={(0.02,0.98)}},
]

\addplot[name path=N8qlo, draw=none, forget plot]
    table[x=N8_bytes, y=N8_lo] {figures/communication_humanoid.dat};
\addplot[name path=N8qhi, draw=none, forget plot]
    table[x=N8_bytes, y=N8_hi] {figures/communication_humanoid.dat};
\addplot[colorD, opacity=0.15, forget plot] fill between[of=N8qlo and N8qhi];
\addplot[colorD, line width=0.9pt, solid]
    table[x=N8_bytes, y=N8_mean] {figures/communication_humanoid.dat};
\addlegendentry{64 bit};

\addplot[name path=q8qlo, draw=none, forget plot]
    table[x=q8_bytes, y=q8_lo] {figures/quant_comm_humanoid.dat};
\addplot[name path=q8qhi, draw=none, forget plot]
    table[x=q8_bytes, y=q8_hi] {figures/quant_comm_humanoid.dat};
\addplot[colorA, opacity=0.15, forget plot] fill between[of=q8qlo and q8qhi];
\addplot[colorA, line width=0.9pt, dashed]
    table[x=q8_bytes, y=q8_mean] {figures/quant_comm_humanoid.dat};
\addlegendentry{8 bit};

\addplot[name path=q6qlo, draw=none, forget plot]
    table[x=q6_bytes, y=q6_lo] {figures/quant_comm_humanoid.dat};
\addplot[name path=q6qhi, draw=none, forget plot]
    table[x=q6_bytes, y=q6_hi] {figures/quant_comm_humanoid.dat};
\addplot[colorB, opacity=0.15, forget plot] fill between[of=q6qlo and q6qhi];
\addplot[colorB, line width=0.9pt, dash dot]
    table[x=q6_bytes, y=q6_mean] {figures/quant_comm_humanoid.dat};
\addlegendentry{6 bit};
\end{axis}
\end{tikzpicture}

    \caption{Humanoid-v4}
  \end{subfigure}
  \caption{Return against cumulative uplink traffic at $N = 8$, per bit width,
  on the same logarithmic axis as Figure~\ref{fig:communication}. The 64-bit
  curve is the $N = 8$ series of that figure. The solid line is the mean over
  five seeds and the band spans the minimum and maximum seed. Quantization
  translates a curve to the left by $\log_{10}(64/b)$ decades, $0.90$ at 8 bits
  and $1.03$ at 6, without changing its shape whenever the bit width is above
  the threshold: the same return is reached with an order of magnitude less
  traffic. Humanoid at 6 bits is the one curve that drops as well as shifts.}
  \label{fig:quantcomm}
\end{figure}

\begin{figure}[t]
  \centering
  \begin{subfigure}[b]{0.49\textwidth}
    %% Final return against bit width, Swimmer-v4, one series per agent count.
%% Mean +- sd over five seeds (same statistic as Table 1); the 64-bit points are
%% the unquantized N=2,4,8 columns of Table 1. Per-agent uplink is 2d*b/8 bytes,
%% so the x axis is proportional to it. Series are offset by a few percent in x
%% so that the error bars do not overlap.
%% Data: figures/quant_tradeoff_swimmer.dat (scripts/export_figures.py).
\begin{tikzpicture}
\begin{axis}[
    width=\linewidth,
    height=0.78\linewidth,
    xlabel={Bits per transmitted value},
    ylabel={Final return},
    xmode=log, log basis x=2,
    xmin=5, xmax=80,
    xtick={6,8,16,32,64},
    xticklabels={6,8,16,32,64},
    ymin=115, ymax=140,
    ytick distance=5,
    ymajorgrids=true, xmajorgrids=true, grid style=dashed,
    tick label style={font=\footnotesize},
    label style={font=\footnotesize},
    legend style={font=\footnotesize, cells={anchor=west},
                  anchor=south, at={(0.5,0.03)}, legend columns=3},
]

\addplot[colorB, line width=0.9pt, mark=*, mark size=2pt,
        error bars/.cd, y dir=both, y explicit]
    table[x expr=\thisrow{bits}*0.93, y=N2_mean, y error=N2_sd]
    {figures/quant_tradeoff_swimmer.dat};
\addlegendentry{$N=2$};

\addplot[colorC, line width=0.9pt, mark=square*, mark size=2pt,
        error bars/.cd, y dir=both, y explicit]
    table[x expr=\thisrow{bits}*1.0, y=N4_mean, y error=N4_sd]
    {figures/quant_tradeoff_swimmer.dat};
\addlegendentry{$N=4$};

\addplot[colorD, line width=0.9pt, mark=triangle*, mark size=2pt,
        error bars/.cd, y dir=both, y explicit]
    table[x expr=\thisrow{bits}*1.075, y=N8_mean, y error=N8_sd]
    {figures/quant_tradeoff_swimmer.dat};
\addlegendentry{$N=8$};
\end{axis}
\end{tikzpicture}

    \caption{Swimmer-v4}
  \end{subfigure}
  \hfill
  \begin{subfigure}[b]{0.49\textwidth}
    %% Final return against bit width, Hopper-v4, one series per agent count.
%% Mean +- sd over five seeds (same statistic as Table 1); the 64-bit points are
%% the unquantized N=2,4,8 columns of Table 1. Per-agent uplink is 2d*b/8 bytes,
%% so the x axis is proportional to it. Series are offset by a few percent in x
%% so that the error bars do not overlap.
%% Data: figures/quant_tradeoff_hopper.dat (scripts/export_figures.py).
\begin{tikzpicture}
\begin{axis}[
    width=\linewidth,
    height=0.78\linewidth,
    xlabel={Bits per transmitted value},
    ylabel={Final return},
    xmode=log, log basis x=2,
    xmin=5, xmax=80,
    xtick={6,8,16,32,64},
    xticklabels={6,8,16,32,64},
    ymin=0, ymax=3600,
    ytick distance=900,
    scaled y ticks=base 10:-3,
    ymajorgrids=true, xmajorgrids=true, grid style=dashed,
    tick label style={font=\footnotesize},
    label style={font=\footnotesize},
    legend style={font=\footnotesize, cells={anchor=west},
                  anchor=north, at={(0.5,0.97)}, legend columns=3},
]

\addplot[colorB, line width=0.9pt, mark=*, mark size=2pt,
        error bars/.cd, y dir=both, y explicit]
    table[x expr=\thisrow{bits}*0.93, y=N2_mean, y error=N2_sd]
    {figures/quant_tradeoff_hopper.dat};
\addlegendentry{$N=2$};

\addplot[colorC, line width=0.9pt, mark=square*, mark size=2pt,
        error bars/.cd, y dir=both, y explicit]
    table[x expr=\thisrow{bits}*1.0, y=N4_mean, y error=N4_sd]
    {figures/quant_tradeoff_hopper.dat};
\addlegendentry{$N=4$};

\addplot[colorD, line width=0.9pt, mark=triangle*, mark size=2pt,
        error bars/.cd, y dir=both, y explicit]
    table[x expr=\thisrow{bits}*1.075, y=N8_mean, y error=N8_sd]
    {figures/quant_tradeoff_hopper.dat};
\addlegendentry{$N=8$};
\end{axis}
\end{tikzpicture}

    \caption{Hopper-v4}
  \end{subfigure}

  \vspace{0.6em}
  \begin{subfigure}[b]{0.49\textwidth}
    %% Final return against bit width, Humanoid-v4, one series per agent count.
%% Mean +- sd over five seeds (same statistic as Table 1); the 64-bit points are
%% the unquantized N=2,4,8 columns of Table 1. Per-agent uplink is 2d*b/8 bytes,
%% so the x axis is proportional to it. Series are offset by a few percent in x
%% so that the error bars do not overlap.
%% Data: figures/quant_tradeoff_humanoid.dat (scripts/export_figures.py).
\begin{tikzpicture}
\begin{axis}[
    width=\linewidth,
    height=0.78\linewidth,
    xlabel={Bits per transmitted value},
    ylabel={Final return},
    xmode=log, log basis x=2,
    xmin=5, xmax=80,
    xtick={6,8,16,32,64},
    xticklabels={6,8,16,32,64},
    ymin=0, ymax=7000,
    ytick distance=1750,
    scaled y ticks=base 10:-3,
    ymajorgrids=true, xmajorgrids=true, grid style=dashed,
    tick label style={font=\footnotesize},
    label style={font=\footnotesize},
    legend style={font=\footnotesize, cells={anchor=west},
                  anchor=south east, at={(0.97,0.03)}, legend columns=1},
]

\addplot[colorB, line width=0.9pt, mark=*, mark size=2pt,
        error bars/.cd, y dir=both, y explicit]
    table[x expr=\thisrow{bits}*0.93, y=N2_mean, y error=N2_sd]
    {figures/quant_tradeoff_humanoid.dat};
\addlegendentry{$N=2$};

\addplot[colorC, line width=0.9pt, mark=square*, mark size=2pt,
        error bars/.cd, y dir=both, y explicit]
    table[x expr=\thisrow{bits}*1.0, y=N4_mean, y error=N4_sd]
    {figures/quant_tradeoff_humanoid.dat};
\addlegendentry{$N=4$};

\addplot[colorD, line width=0.9pt, mark=triangle*, mark size=2pt,
        error bars/.cd, y dir=both, y explicit]
    table[x expr=\thisrow{bits}*1.075, y=N8_mean, y error=N8_sd]
    {figures/quant_tradeoff_humanoid.dat};
\addlegendentry{$N=8$};
\end{axis}
\end{tikzpicture}

    \caption{Humanoid-v4}
  \end{subfigure}
  \caption{Final return against bit width for each agent count, mean $\pm$
  standard deviation across five seeds (Table~\ref{tab:quant}). Per-agent uplink
  is $2db/8$ bytes, so the horizontal axis is proportional to it, and moving
  from the right-hand points to the left-hand ones divides the traffic by $8$
  and $10.7$. Series are offset slightly in $x$ so the error bars stay
  readable. The Swimmer axis is zoomed to the range of the data. A flat line
  is the desired outcome; only Humanoid at 6 bits breaks it.}
  \label{fig:quanttradeoff}
\end{figure}

\paragraph{8 bits cost nothing measurable.} At 8 bits the uplink falls by
$8\times$ on every environment, and Table~\ref{tab:quant} shows no cell in
which the reward pays for it. On Swimmer the quantized means sit within $0.2$
of the unquantized ones at all three agent counts, with a seed spread that is,
if anything, tighter ($\pm 0.4$ against $\pm 0.8$ at $N=8$). On Humanoid, the
environment with the most at stake, the gap is $99$ at $N=4$ and $332$ at
$N=8$, each inside one standard deviation of the reference, and the $N=2$ gap
of $435$ is inside the $\pm 796$ spread of the unquantized runs. On Hopper the
8-bit means are higher than the reference at $N = 4$ and $8$ and lower at
$N = 2$, all within the seed spread. Figure~\ref{fig:quanttradeoff} draws the
same result as flat lines, and Figure~\ref{fig:quantcomm} shows it along the
whole run rather than only at the end: the 8-bit curve has the shape of the
64-bit one, displaced by almost a decade to the left.

In absolute terms the saving is large. A full $N=8$ Humanoid run uplinks
\SI{279}{\giga\byte} at 64 bits and \SI{34.9}{\giga\byte} at 8, so 8-bit
quantization removes \SI{244}{\giga\byte} of agent-to-server traffic from a
single training run at no measurable cost in return. The corresponding figures
are \SI{4.72}{\giga\byte} to \SI{0.59}{\giga\byte} on Hopper and
\SI{623}{\mega\byte} to \SI{78}{\mega\byte} on Swimmer.

\paragraph{The saving compounds with the $\Oh(d)$ one.} Standard-average FedNPG
uplinks $d^2 + d$ values per agent per round at 64 bits; 8-bit FedNPG-ADMM
uplinks $2d$ values at 8 bits. The ratio is $4(d+1)$, which turns the
$2434\times$, $9219\times$ and $363537\times$ reductions of
Section~\ref{sec:baselines} into $1.9 \times 10^4$, $7.4 \times 10^4$ and
$2.9 \times 10^6$ on Swimmer, Hopper and Humanoid. Quantization also changes
the ordering in Table~\ref{tab:measuredcomm}. At 64 bits FedNPG-ADMM transmits
twice what FedPPO does; at 8 bits it transmits $16d$ bits per agent per
iteration against FedPPO's $64d$, a quarter of the unquantized first-order
baseline. FedPPO's parameter delta could in principle be quantized as well,
which we did not test, so this is a statement about the configurations we ran
and not a ranking of the two methods under equal compression.

The converse does not hold: quantization cannot rescue standard-average
FedNPG. To match 8-bit FedNPG-ADMM, its $d^2 + d$ values would have to travel
at $16/(d+1)$ bits each, $3.3 \times 10^{-3}$ bits on Swimmer and less on the
larger policies. No quantizer reaches below one bit per value, so the
$\Oh(d^2) \to \Oh(d)$ reformulation is a precondition for the precision axis
to matter, not an alternative to it.

\paragraph{6 bits: a threshold that depends on $d$.} At 6 bits the saving
rises to $10.7\times$. Swimmer ($d = \num{4868}$) and Hopper
($d = \num{18438}$) keep their returns, while Humanoid ($d = \num{727074}$)
collapses at every agent count, to between $13\%$ and $24\%$ of the
unquantized return. Two qualifications apply to the surviving environments. On
Swimmer the 6-bit spread widens at $N = 2$ and $N = 4$ ($\pm 7.4$ and $\pm 3.4$
against $\pm 1.3$ and $\pm 0.8$). The $N=2$ mean of $128.9$ is pulled up by a
single seed at $142.0$, with the other four between $124.9$ and $126.8$, and at
$N=4$ two seeds end near $123$; this is the first visible effect of the
quantization noise, and it is gone at $N = 8$. On Hopper the 6-bit means show no
ordering in bit width ($1573$, $1853$, $1837$ against $2140$, $1658$, $1605$),
and the differences are inside the seed spread that Section~\ref{sec:hopper}
documents for this environment. The Hopper column measures seed variance, not
quantization.

That the threshold moves with $d$ is consistent with how the quantizer sets
its step. The step $s$ is fixed by the range of the whole tensor, and in a
vector of $\num{727074}$ entries the extremes sit much further out in the
tails than in one of $\num{4868}$. Most entries are then small compared with
$s$, and what reaches the server for them is dominated by rounding noise,
correct only in expectation. We have not measured the per-tensor ranges, so this reading is a
hypothesis; a per-bit sweep on a fourth environment of intermediate $d$ would
test it.

The collapsed Humanoid runs are at least consistent with the reason we
chose an unbiased quantizer. At 6 bits, return grows $3.0\times$ from $N=2$ to
$N=8$ ($497 \to 1514$), against $1.9\times$ for the unquantized runs
($3391 \to 6329$), so the fraction of the reference retained rises from
$15\%$ and $13\%$ at $N = 2$ and $4$ to $24\%$ at $N=8$. Each agent's
quantization error is independent and zero-mean, so the server's average over
$N$ agents divides its variance by $N$, and more agents buy back part of the
fidelity the coarse quantizer took away. A biased quantizer would leave the
same systematic error after averaging. With three points and a non-monotone
middle one, this is suggestive rather than conclusive.

\subsection{What the byte counts omit}

The reported bytes are $d \cdot b / 8$ per transmission and exclude the
quantization metadata a real deployment would also send: $\min x$ and the
scale, two \texttt{float32} values or 64 bits per transmission. Relative to a
$d \cdot b$-bit payload the overhead is $64/(d\,b)$. At the bit widths we run
this is at most $0.22\%$ (Swimmer, 6 bits), $0.06\%$ on Hopper and below
$0.002\%$ on Humanoid, so the $8\times$ saving is $7.99\times$ once the
metadata is paid for. The downlink broadcast of $\bthe$ and $\by$ is not
counted either, for the same reason it is not counted in the paper. Neither
omission changes the ordering in Table~\ref{tab:quant}, but both make the
reported compression slightly optimistic.

\paragraph{Skipped steps.} Quantization noise can in principle drive
$(\sum_i \bg_i)^{\top}\by$ negative, leaving line~12's step size undefined
(Section~\ref{sec:implementation}). On Hopper and Humanoid this never occurs
in any of the 60 quantized runs. On Swimmer it does, but it does not track the
bit width: the fifteen 8-bit runs skip 51 of their \num{15000} iterations in
total (at most 12 in one run), while the fifteen 6-bit runs skip 3. Coarser
quantization does not make the step more often undefined.
